\chapter{Testing}
\label{chap:testing}

Seal is tested at two levels: \textbf{unit tests} co-located in \verb|src/|
modules, and \textbf{integration tests} in \verb|tests/| that drive the real
server over HTTP and WebSocket. At the time of writing the \verb|tests/| tree
holds 72 \verb|#[tokio::test]| integration cases, plus unit-test modules in seven
source files (\verb|auth|, \verb|config|, \verb|db_ops|, \verb|error|,
\verb|rate_limit|, \verb|validate|, \verb|ws|).

\begin{quote}
\textbf{Doc drift note.} The \verb|README| still advertises ``62 native Rust
integration tests''; the actual count has since grown. Treat the source tree, not
the prose, as authoritative.
\end{quote}

\section{The Integration Harness}

\verb|tests/common/mod.rs| defines \verb|TestServer|, which spins up the genuine
application stack against a throwaway database:

\begin{lstlisting}
let tempdir = tempfile::tempdir()?;                 // isolated, auto-cleaned
let db_path = tempdir.path().join("test.lance");
let cfg = Config::for_test(db_path, "test-secret-key-for-tests".into())?;
let conn = db::connect(&cfg.database_path).await?;
db::init_db(&conn).await?;
let state = AppState { cfg: Arc::new(cfg), conn,
                       rate_limiter: ..., ws_connections: ... };
let router = build_router(state.clone())
    .into_make_service_with_connect_info::<SocketAddr>();
let listener = TcpListener::bind("127.0.0.1:0").await?;  // random free port
\end{lstlisting}

The server runs in a spawned task with graceful shutdown wired to a
\verb|oneshot| channel held in the struct; when the \verb|TestServer| drops, the
server stops and the \verb|TempDir| is deleted. Two properties make the suite
robust and fast:

\begin{itemize}
  \item \textbf{Real router, real DB.} Tests exercise \verb|build_router| and
        LanceDB exactly as production does --- no mocks, no in-memory
        substitutes. Only the config source (\verb|for_test|) and the storage
        location (a temp dir) differ.
  \item \textbf{Parallel-safe.} Each test gets its own port, its own database,
        and a config that ignores process env vars, so the default parallel test
        runner has nothing to race on.
\end{itemize}

\verb|TestServer| also exposes \verb|base_url|, a \verb|client()| builder
(\verb|reqwest|), the \verb|state| (so a test can seed rows directly through
\verb|db_ops|), and a \verb|url(path)| helper.

\section{The Testing Philosophy: Seed via Internals, Assert via API}

The suite's guiding pattern is to \emph{set up} state through the library
internals (writing rows with \verb|db_ops::append|, minting tokens with
\verb|create_token|) and then to \emph{exercise and assert} through the public
HTTP/WS surface. This keeps tests fast and precise while still covering the real
request path. Each test file keeps its own small local helpers (\verb|register|,
\verb|create_channel|, \verb|open_ws|, \dots) rather than a shared ``kitchen
sink'', so a test reads top-to-bottom.

\section{What the Integration Files Cover}

\begin{center}
\begin{longtable}{@{}lp{0.60\textwidth}@{}}
\toprule
\textbf{File} & \textbf{Focus} \\
\midrule
\endhead
\verb|tests/auth.rs| & Registration, login, duplicate-user and bad-credential paths, JWT acceptance/rejection, bcrypt compatibility, and rate limiting. \\
\verb|tests/users.rs| & DM-contact derivation, prefix search (self-exclusion, the 20 cap), and public-key retrieval. \\
\verb|tests/channels.rs| & Channel create/list/browse/join, member invite, membership authorization, duplicate-name and duplicate-member rejection, and member-key listing. \\
\verb|tests/messages.rs| & DM history merging and \verb|after| filtering, channel history scoping, the REST channel-send path, and attachments. \\
\verb|tests/websocket.rs| & Upgrade and token gating, DM relay (including self-copy and multi-device echo), channel fan-out, and non-member rejection. \\
\verb|tests/migration.rs| & The messages-table column migration and its idempotency on re-run. \\
\verb|tests/static_assets.rs| & Serving the embedded \verb|index.html| and static assets with correct content types. \\
\bottomrule
\end{longtable}
\end{center}

\section{Unit Tests}

The co-located \verb|#[cfg(test)]| modules test pure logic in isolation:
\verb|config| (the cascade and env precedence), \verb|auth| (hashing, tokens,
expiry), \verb|db_ops| (row builders, null/missing/wrong-type handling, the
message materializer), \verb|error| (status/body mapping and internal masking),
\verb|rate_limit| (boundary, per-IP, reset), \verb|validate| (accepted/rejected
inputs), and \verb|ws| (registry fan-out and cleanup).

\section{Running the Tests}

\begin{lstlisting}
make test          # cargo test  -- everything
make test-quiet    # cargo test --quiet
cargo test --test websocket        # one integration file
cargo test config::                # a module's unit tests
\end{lstlisting}

\section{Coverage and Intentionally-Untested Paths}

Line coverage is measured with \verb|cargo llvm-cov| (see the coverage CI
workflow in Chapter~\ref{chap:build}). A few paths are deliberately left
uncovered because exercising them would add machinery out of proportion to their
risk: \verb|src/main.rs| (the thin bootstrap, covered in spirit by the harness),
the database-failure error closures inside \verb|db_ops| (they require injecting a
broken connection), and some embedded-asset ``file missing / not UTF-8'' branches
in \verb|lib.rs| that cannot occur with the assets that are actually compiled in.
These omissions are documented rather than papered over.
